\documentclass{scrreprt}

\usepackage{aligned-overset}
\usepackage{amsmath}
\usepackage{amsthm}
\usepackage{amssymb}
\usepackage{bm}
\usepackage[inline,shortlabels]{enumitem}
\usepackage{hyperref}
\usepackage[utf8]{inputenc}
\usepackage{multicol}
\usepackage{mathtools}
\usepackage{pdflscape}
\usepackage{physics}
\usepackage{tabularx}
\usepackage[table]{xcolor}
\usepackage{titling}
\usepackage{fancyhdr}
\usepackage{xfrac}
\usepackage{pgfplots}

\pgfplotsset{compat = newest}
\usepgfplotslibrary{fillbetween}
\usetikzlibrary{calc}


\author{Karsten Lehmann}
\date{WiSe 2025/26}
\title{Übungsblatt 01\\INF-B-270, Formale Systeme}

\setlength{\parindent}{0pt}

\setlength{\headheight}{26pt}
\pagestyle{fancy}
\fancyhf{}
\lhead{\thetitle}
\rhead{\theauthor}
\lfoot{\thedate}
\rfoot{Seite \thepage}

\begin{document}
\paragraph{Aufgabe 1}
Gegeben sind ein beliebiges Alphabet $\Sigma$ und die Sprachen
$L_1, L_2, L_3, L_4 \subseteq \Sigma^*$.
Zeigen Sie, dass die Operationen $\cup, \circ$ und $^*$ monoton sind, d.h.
für $L_1 \subseteq L_3$ und $L_2 \subseteq L_4$ gilt:
\begin{enumerate}[(a)]
\item $L_1 \cup L_2 \subseteq L_3 \cup L_4$

  \subparagraph{Lsg.} Sei $u \in L_1 \cup L_2$ beliebig.
  Dann gilt nach der Definition der Vereinigung, dass
  \begin{enumerate}[(1)]
  \item $u \in L_1$, dann da $L_1 \subseteq L_3$ auch $u \in L_3$ und damit
    auch $u \in L_3 \cup L_4$
  \item $u \in L_2$, dann da $L_2 \subseteq L_4$ auch $u \in L_4$ und damit
    auch $u \in L_3 \cup L_4$
  \end{enumerate}

\item $L_1 \circ L_2 \subseteq L_3 \circ L_4$

  \subparagraph{Lsg.} Sei $u = wv \in L_1 \circ L_2$.
  Dann sind $w \in L_1$ und $v \in L_2$.
  Da $L_1 \subseteq L_3$ und $L_2 \subseteq L_4$ gilt somit auch
  $w \in L_3$ und $v \in L_4$ und damit schließlich auch
  $wv = u \in L_3 \cup L_4$.
\item $L_1^* \subseteq L_3^*$

  \subparagraph{Lsg.} Beweis über vollständige Induktion.in zwei Schritten:

  \textbf{Behauptung:}
  \[
    P\qty\big(n) \colon L_1^n  \subseteq L_3^n
  \]

  \textbf{Induktionsanfang:}
  \begin{flalign*}
    P\qty\big(0) &\colon L_1^0 = \qty\big{\epsilon} \subseteq \qty\big{\epsilon} = L_3^0 \\
    P\qty\big(1) &\colon L_1 \subseteq  L_3
  \end{flalign*}

  \textbf{Induktionsschritt:} Angenommen es wäre nun $P\qty\big(n)$ für ein
  beliebiges $n \in \mathbb{N}$ wahr (Induktionsvoraussetzung (IV)), dann
  \begin{flalign*}
    P\qty\big(n + 1) \colon L_1^{n+1} \subseteq L_3^{n + 1}
    &\iff L_1 \circ L_1^n \subseteq L_3 \circ L_3^n
  \end{flalign*}
  Sei nun $u = wv \in L_1 \circ L_1^n$ beliebig mit $w \in L_1$ und $v = L_1^n$,
  dann ist auch $w \in L_3$ und nach der Induktionsvoraussetzung $v \in L_3^n$.

  Somit ist $u \in L_3 \circ L_3^n$ und da $u$ beliebig, die Behauptung
  $P\qty\big(n + 1)$ wahr.

  Und aus dem Satz über die vollständige Induktion folgt die Behauptung.

  Analog lässt sich auch über vollständige Induktion zeigen, dass
  \[
    L_1^0 \cup L_1^1 \cup \ldots L_1^n \subseteq L_3^0 \cup L_3^1 \cup \ldots L_3^n
  \]

\end{enumerate}

\paragraph{Aufgabe 2}
Zeigen oder widerlegen Sie die folgenden Identitäten.
Dabei sind $L, L_1, L_2, L_3$ beliebige Sprachen.

\begin{enumerate}[(a)]
\item $L_1 \circ \qty(L_2 \cup L_3) = L_1 \circ L_2 \cup L_1 \circ L_3$,
  $L_1 \circ \qty(L_2 \cup L_3) = \qty(L_1 \circ L_2+) \cap \qty(L_1 \circ L_3)$.

  \subparagraph{Lsg.} Die Distributivität über die Vereinigung wurde bereits in
  der ersten Vorlesung gezeigt.
  Für den Schnitt hingegen gilt die Distributivität hingegen nicht, denn seien
  $L_1 = \qty\big{\text{a}, \text{aa}}, L_2 = \qty\big{\text{b}}$
  und $L_3 = \qty\big{\text{ab}}$, dann ist
  \begin{flalign*}
    L_1 \circ \qty(L_2 \cap L_3)
    &= \qty\big{\text{a}, \text{aa}} \circ \qty(\qty\big{\text{b}} \cap \qty\big{\text{ab}}) \\
    &= \qty\big{\text{a}, \text{aa}} \circ \emptyset \\
    &= \emptyset \\
    &\ne \qty\big{\text{aab}} \\
    &= \qty\big{\text{ab}, \text{aab}} \cap \qty\big{\text{aab}, \text{aaab}} \\
    &= \qty\big(\qty\big{\text{a}, \text{aa}} \circ \qty(\qty\big{\text{b}})) \cap \qty\big(\qty\big{\text{a}, \text{aa}} \circ \qty(\qty\big{\text{ab}})) \\
    &= \qty\big(L_1 \circ L_2) \cap \qty\big(L_1 \circ L_3)
  \end{flalign*}

\setcounter{enumi}{3}
\item $\emptyset^* = \qty\big{\epsilon}$

  \subparagraph{Lsg.} Das ist korrekt, da per Definition
  $\emptyset^0 = \qty\big{\epsilon}$ und $\emptyset^n = \emptyset$ für alle $n \ge 1$.
  Somit ist
  $\emptyset^0 \cup \emptyset^1 \cup \ldots \cup \emptyset^n = \qty\big{\epsilon}$.
\end{enumerate}

\paragraph{Aufgabe 3}
Gegeben ist die Grammatik $G = \langle V, \Sigma, P, S \rangle$
mit $V = \qty\big{S, B}, \Sigma = \qty\big{a, b, c}$ und
$P = \qty\big{S \to aSBc, S \to abc, cB \to B, cB \to Bc, bB \to bb}$.

\begin{enumerate}[(a)]
\item Von welchem maximalen Typ ist $G$?
  Begründen Sie Ihre Antwort.

  \subparagraph{Lsg.} Die Sprache ist \textbf{nicht} kontextfrei, da in den
  Regeln auf der linken Seite zum Teil mehr als nur eine Variable vorkommt.

  Weiter ist die Sprache auch \textbf{nicht} kontextsensitiv, da für
  $cB \to B$ nicht $\abs{cB} \le \abs{B}$.

  Somit handelt es sich um eine Grammatik vom Typ 0.

\item Bestimmen Sie drei Wörter $w$ der Sprache $L\qty\big(G)$.
  Begründen Sie Ihre Antwort.

  \subparagraph{Lsg.} Zum Beispiel
  \begin{enumerate}[1)]
  \item $S \to abc$
  \item $S \to aSBc \to aabcBc \to aabBcc \to aabbcc$
  \item $S \to aSBc \to aabcBc \to aabBc \to aabbc$
  \end{enumerate}

\newpage
\item Gilt $\epsilon \in L\qty\big(G)$?
  Begründen Sie Ihre Antwort!

  \subparagraph{Lsg.} Nein, da egal ob $S$ nach der ersten oder
  zweiten Regel aufgelöst wird, es steht immer ein $a$ am Anfang, welches auch
  mit keiner weiteren Regel mehr entfernt werden kann.

\item Beschreiben Sie die durch $G$ erzeugte Sprache $L\qty\big(G)$.

  \subparagraph{Lsg.} In allen Regeln werden gleich viele $a$'s und $b$'s
  hinzugefügt und nicht wieder entfernt.

  Zusätzlich werden gleich viele $c$'s wie $a$'s und $b$'s eingefügt, allerdings
  eventuell auch durch die 3. Regel wieder entfernt.

  Somit ist $L\qty\big(G) = \qty{a^nb^nc^m \;{\big |}\; 1 \leq m \leq n}$
\end{enumerate}
\end{document}
